\section{Blind transport and singular acquisition}\label{sec:singular}
The exact criterion also applies when every displayed one-step stability coefficient is one, and when no affine density exists. This gives a sharp law different from the smoothing law of Section~\ref{sec:geometry}.

Let $X$ have a specified prepared law $\mu$ on a compact subset $E$ of a finite-dimensional Euclidean space. A protocol consists of one paid capture and $n-1$ paid blind calls. Only the action \emph{capture} is legal initially; it reports $X$ exactly. At subsequent nonterminal phases only \emph{advance} is legal and its report is a fixed symbol. The physical variable $X$ remains unchanged. After exactly $n$ calls the observer predicts $X$ under squared loss. The terminal scored outcome records $X$ after the decision, so bounded continuous functions of that outcome are executable terminal tests. It is unavailable before the prediction; an additional physical audit call, when needed, is counted separately. Thus the capture/advance interface is part of the experiment, not an uncounted clock supplied to the observer.

Define the actual-law quantization error
\[
 q_L(\mu)=\inf_{c_1,\ldots,c_L\in\operatorname{conv}(E)}
               \int\min_j\|x-c_j\|_2^2\mu(dx),\qquad L\ge1.
\]
The centers may be restricted to the compact convex hull, so a minimizer exists. The countable future tests on $E$ give the compact predictive carrier $\mathcal P(E)$; the capture law on this carrier is the pushforward of $\mu$ by $x\mapsto\delta_x$. Blind updates are the identity. Affine mixture closure holds for both raw kernels. No ambient smooth parameterization of $\mu$ is used.

\begin{theorem}[Exact blind-transport law]\label{thm:blind}
For this experiment, every external width profile $\mathbf m=(m_1,\ldots,m_n)$ satisfies
\begin{equation}\label{eq:blind_external}
 R_{\ext}(\mathbf m)=q_{\min_t m_t}(\mu).
\end{equation}
The stationary internally exact-$n$ class is nonempty precisely when $M\ge n+1$, and in that range
\begin{equation}\label{eq:blind_internal}
 R_{\aut}(n,M)=q_{\lfloor(M-1)/n\rfloor}(\mu).
\end{equation}
In particular, these conclusions hold for singular continuous and countably atomic acquired laws. They are assertions for the declared protocol and task, not for an experiment permitting repeated captures.
\end{theorem}
\begin{proof}
Fix a cut $t\ge1$. All future reports are constant, and fresh transition randomness is independent of $X$ conditional on the current label $Z_t$. Hence the terminal prediction $U$ satisfies the Markov relation $X-Z_t-U$. Conditional squared-loss decomposition gives
\[
 \E\|X-U\|^2\ge \E\|X-\E[X\mid Z_t]\|^2\ge q_{|\supp Z_t|}(\mu).
\]
The last inequality remains valid for randomized encoders: the conditional means give at most $|\supp Z_t|$ centers, and pointwise least squared distance is no larger than the encoder's conditional average. This lower holds at every cut. Conversely, quantize $X$ once using $L=\min_t m_t$ centers and carry the resulting label unchanged through all blind calls. The external clock supplies the legal phase and terminal request. This proves \eqref{eq:blind_external}.

For the autonomous machine, its initial label cannot occur later with positive probability. Its action must be capture initially and advance at every later nonterminal phase; the same stationary action kernel cannot do both. Nor can a nonterminal label be a stopping label. Between cuts $1$ and $n$, the reports are identical and the label dynamics are stationary. If one positive-probability label occurred at two different such cuts, its law of future stopping time would be the same from both, contradicting exact stopping at $n$. Thus these $n$ cut supports and the initial label are disjoint. Their sizes $L_t\ge1$ satisfy $1+\sum_{t=1}^nL_t\le M$, so one is at most $\lfloor(M-1)/n\rfloor$. The cut lower just proved gives \eqref{eq:blind_internal}. To attain it, use one initial capture label and $n$ disjoint copies of an optimal $L$-point quantizer label. The first copy is selected by the capture report, every blind transition advances its phase copy without changing its quantizer index, and the last copy stops and emits the center. The total number of states is $1+nL\le M$. The same construction with $L=1$ proves nonemptiness at $n+1$, and the disjoint-support argument proves necessity.
\end{proof}

This theorem is also an explicit solution of the allocation frontier in Theorem~\ref{thm:completion}. At capture, the law of conditional distributions is a barycentric contraction of the actual law of $\delta_X$. At every later cut a further contraction can only lose information. Carrying the chosen conditional distribution without further pooling makes all later Jensen defects zero, although the blind update has no strict contraction at all. Thus nonsummable bounds on generic Lipschitz sensitivities do not preclude exact causal completion. Conversely, the narrowest retained cut gives a lower unrelated to an artificial terminal density.

\begin{corollary}[Fractal and accumulating atomic laws]\label{cor:singular_rates}
\textup{(i)} Let $\mu$ be the equal-weight two-map Cantor law with contraction $r\in(0,1/2)$ and set $s=\log2/\log(1/r)$. Then
\[
 q_L(\mu)\asymp_r L^{-2/s}.
\]
\textup{(ii)} Let $\alpha>0$, $x_j=2^{-j}$, and
\[
 \mu\{x_j\}=(1-2^{-\alpha})2^{-\alpha(j-1)},\qquad j\ge1.
\]
On the compact support $\{0\}\cup\{x_j:j\ge1\}$ one has
\[
 q_L(\mu)\asymp_\alpha 2^{-(\alpha+2)L}.
\]
Substitution into \eqref{eq:blind_external}--\eqref{eq:blind_internal} gives matching causal laws, with the actual prepared masses unchanged.
\end{corollary}
\begin{proof}
For (i), a level-$k$ cylinder has diameter $r^k$ and mass $2^{-k}$. With $k=\lfloor\log_2L\rfloor$, its $2^k$ cylinder representatives give $q_L\le r^{2k}\le C_rL^{-2/s}$. The gaps separating neighboring cylinders at their first distinct level imply
$\mu(B(x,\rho))\le C_r\rho^s$ for every $x$ and $0<\rho\le1$: choose $k$ with $r^{k+1}<\rho\le r^k$; an interval of length $2\rho$ meets only a bounded number, depending on $r$, of level-$k$ cylinders. For $L$ centers choose $\rho=(2C_rL)^{-1/s}$, after increasing the constant for the finitely many coarse radii. At least half the acquired mass is outside their balls. This gives $q_L\ge\rho^2/2$.

For (ii), retain the first $L-1$ atoms exactly and place the remaining center at zero. Then
\[
 q_L\le\sum_{j\ge L}(1-2^{-\alpha})2^{-\alpha(j-1)}2^{-2j}
       \le C_\alpha 2^{-(\alpha+2)L}.
\]
The $L+1$ intervals centered at $x_j$ with radii $x_j/8$, $1\le j\le L+1$, are pairwise disjoint. At least one contains no one of $L$ proposed centers. Its atom alone contributes at least
$c_\alpha2^{-(\alpha+2)(L+1)}$ to the squared error. This lower is uniform over centers and proves the claim.
\end{proof}
The support in (ii) accumulates, and its effective finite-resolution curve is exponential rather than the power predicted by an ambient one-dimensional chart. Its proof uses the actual atom masses. The result does not assert a corresponding formula for arbitrary singular laws: for those, Theorem~\ref{thm:blind} retains the exact quantity $q_L(\mu)$.

In Theorem~\ref{thm:blind}, $m_n$ counts terminal stopping/readout labels and the decision may be issued immediately upon entering one. The post-decision audit cannot break the cut Markov relation $X-Z_t-U$. In the accumulating-atom lower bound, absence of a center from the selected interval means that every center is at distance at least $x_j/8$ from its atom; this gives exactly the asserted single-atom contribution.
